\section{Conceptual framework}\label{sec:framework}

\subsection{Four motives and their fingerprints}\label{sec:motives}

Observed stablecoin flows on Thai rails are the sum of four motives. \emph{Trading and offshore bridging}: residents buy
USDT as collateral or as the vehicle to reach offshore derivatives. \emph{Dollar saving}: residents hold USDT as a
store of value. \emph{Cross-border transfer}: migrant workers, traders and tourists move value across borders.
\emph{Illicit pass-through}: criminal proceeds enter as baht and leave as USDT. Each motive leaves a distinct
fingerprint, summarised in \cref{tab:motives}, and the identification strategy is to remove the first three before
interpreting the residual (\cref{fig:funnel}).

\begin{table}[!htbp]
\centering\small
\caption{Motives for holding and moving stablecoins on Thai rails, their fingerprints and how they are tested}
\label{tab:motives}
\begin{tabularx}{\textwidth}{L{2.6cm}XXc}
\toprule
Motive & Fingerprint in the data & Main confounder & Hypothesis \\
\midrule
Trading and offshore bridge & USDT volume co-moves with crypto volatility and offshore funding rates; withdrawals go
to labelled exchange hot wallets and return after gains & Arbitrageurs move as fast as mules & H1, H2 \\
\addlinespace
Dollar saving & Long holding periods; USDT/THB premium over the official rate rises when the baht weakens & Little
in Thailand (low inflation) & H1 \\
\addlinespace
Cross-border transfer & Small, regular amounts; peaks after paydays and before festivals; destinations are P2P desks in
neighbouring countries & Tainted P2P merchants & H1, H6 \\
\addlinespace
Illicit pass-through & Young accounts funded by many bank accounts; deposits structured below 50,000 baht; exits within
minutes; forward paths to flagged addresses; responds to crackdowns, not prices & Honest P2P dealers, camouflaged mules
& H3--H6 \\
\bottomrule
\end{tabularx}
\end{table}

\begin{figure}[!htbp]
\centering
\begin{tikzpicture}[
  font=\sffamily\scriptsize, >={Stealth[length=2.2mm]}, every node/.append style={execute at begin node={\hyphenpenalty=10000\exhyphenpenalty=10000}},
  stage/.style={draw=faint, line width=0.5pt, fill=rule!45, text width=2.5cm, align=center, inner sep=4pt},
  last/.style={draw=thai, line width=1pt, fill=thai!10, text width=2.5cm, align=center, inner sep=4pt},
  note/.style={text width=2.75cm, align=center, text=muted, font=\sffamily\tiny, inner sep=1pt},
  test/.style={text width=2.95cm, align=center, text=thai, font=\sffamily\tiny\bfseries, inner sep=1pt}]
\node[stage, minimum height=3.3cm] (a) at (0,0) {\textbf{All stablecoin flows on Thai rails}\\[3pt]exchange purchases, withdrawals,\\ P2P and OTC trades};
\node[stage, minimum height=2.7cm, right=0.55cm of a] (b) {\textbf{Net of trading and offshore bridging}};
\node[stage, minimum height=2.1cm, right=0.55cm of b] (c) {\textbf{Net of dollar saving}};
\node[stage, minimum height=1.6cm, right=0.55cm of c] (d) {\textbf{Net of cross-border transfers}};
\node[last, minimum height=1.1cm, right=0.55cm of d] (e) {\textbf{Residual to explain}};
\foreach \u/\v in {a/b,b/c,c/d,d/e} \draw[->, draw=muted, line width=0.6pt] (\u.east) -- (\v.west);
\node[note, below=0.25cm of a.south] (na) {observed volume, withdrawals, premium, network mix};
\node[note, anchor=north] at (na.north -| b) {control asset on the same platform; crypto volatility; exchange-labelled destinations};
\node[note, anchor=north] at (na.north -| c) {USDT/THB premium versus the reference rate; holding periods};
\node[note, anchor=north] at (na.north -| d) {payday and festival seasonality; destinations in neighbouring countries};
\node[note, anchor=north] at (na.north -| e) {structuring, pass-through speed, fan-in, forward exposure};
\node[test, above=0.2cm of a.north] (ta) {H1 excess demand};
\node[test, anchor=south] at (ta.south -| b) {H2 USDT vs.\ other assets};
\node[test, anchor=south] at (ta.south -| c) {H1 premium};
\node[test, anchor=south] at (ta.south -| d) {H1, H6 destinations};
\node[test, anchor=south] at (ta.south -| e) {H3 dwell \textperiodcentered{} H4 bunching\\H5 shocks \textperiodcentered{} H6 exposure};
\end{tikzpicture}
\caption{Identification strategy. Observed flows are the sum of four motives. Each legitimate motive is removed using
its fingerprint (grey notes) before the residual is interpreted; red labels give the pre-registered hypotheses of
\cref{sec:plan} that test each step. Box heights are schematic, not estimates.}
\label{fig:funnel}
\end{figure}

\subsection{Structuring at a verification notch}\label{sec:bunching}

Consider deposits of baht onto a licensed platform. Each deposit has an intended amount $z$ and funds an asset class
$a\in\{S,O\}$, where $S$ denotes USDT and $O$ any other crypto-asset. Banks require face verification for transfers of
at least $T$ baht. Let $h_a(z)$ be the counterfactual density of intended amounts of honest deposits funding $a$, and
write $H_a(A)=\int_A h_a(z)\,dz$ for the counterfactual mass in a set $A$.

Honest depositors whose intended amount lies in a response window $[T,T+R)$ relocate to a bunching window $[T-\delta,T)$
with probability $\pi_a(z)$; they find the scan a hassle, or simply prefer to split the transfer. Launderers add a mass
$I_a$ of \emph{structured} deposits placed in $[T-\delta,T)$ to avoid the verification step that would expose a mule's
identity. The observed excess mass below the threshold and the responding share are
\begin{equation}\label{eq:excess}
B_a=\int_T^{T+R}\pi_a(z)\,h_a(z)\,dz + I_a,
\qquad
\beta_a=\frac{B_a}{H_a\!\left([T,T+R)\right)}=\bar\pi_a+\frac{I_a}{H_a\!\left([T,T+R)\right)},
\end{equation}
where $\bar\pi_a$ is the average relocation probability over the response window.

\begin{assumption}[Asset-invariant hassle avoidance]\label{as:a1}
$\bar\pi_S=\bar\pi_O$: the share of honest deposits intended for the response window that relocate below $T$ does
not depend on which asset the deposit funds.
\end{assumption}
\begin{assumption}[Channel specificity]\label{as:a2}
$I_O=0$: structured deposits fund the stablecoin channel, which offers immediate, liquid exit.
\end{assumption}
\begin{assumption}[Smooth counterfactual]\label{as:a3}
$h_a$ is smooth in $\log z$ on a neighbourhood of $[T-\delta,T+R)$, so that it is identified from the observed density
outside that window, up to multiplicative round-number effects.
\end{assumption}

\begin{proposition}[Difference in bunching]\label{prop:bunching}
Under \cref{as:a1,as:a2,as:a3},
\[
\Delta\beta\equiv\beta_S-\beta_O=\frac{I_S}{H_S\!\left([T,T+R)\right)}\geq 0,
\]
with equality if and only if $I_S=0$. The structured stablecoin volume is identified as
$I_S=\Delta\beta\cdot H_S([T,T+R))$.
\end{proposition}

The proof is immediate from \cref{eq:excess} and is given in \cref{app:proofs}. Two features matter in practice. First,
the normalisation: the classic statistic divides the excess by the counterfactual density just below $T$. When the two
asset classes have amount distributions of different shapes---stablecoin deposits are larger---that statistic differs
across assets even when $\bar\pi_S=\bar\pi_O$. In the simulation below, a test based on it rejected a true null in 14\%
of draws at a nominal 5\% level; the responding share $\beta_a$ in \cref{eq:excess} reduced this to 7.5\%. Second,
\cref{as:a1} can be relaxed. If deposits observed before the verification rule exist, the change in $\Delta\beta$ across
the rule's introduction identifies $I_S$ under the weaker assumption that asset-specific differences in honest bunching
are stable over time.

\begin{corollary}[Difference-in-difference-in-bunching]\label{cor:didib}
If $\bar\pi_S-\bar\pi_O$ is the same before and after the rule's introduction and no structuring occurs at $T$ before
it, then $\Delta\beta^{\text{post}}-\Delta\beta^{\text{pre}}=I_S^{\text{post}}/H_S^{\text{post}}([T,T+R))$.
\end{corollary}

A placebo version of the test at amounts where no verification step exists ($T'$ = 30,000 or 70,000 baht) provides a
specification check: under correct specification $\Delta\beta(T')$ should be zero.

\subsection{Pass-through and the limits of speed}\label{sec:speed}

A mule account exists to move value, so the time between a baht deposit and the withdrawal of the USDT it bought---its
\emph{dwell time}---is short. But speed is shared with honest populations: arbitrageurs exploit price gaps within
minutes, and peer-to-peer dealers recycle inventory quickly. The following bound formalises why a speed-only rule
cannot achieve high precision when such populations exist.

\begin{lemma}[Speed is not sufficient]\label{lem:speed}
Let $n_M$ mule accounts and $n_F$ honest ``fast'' accounts have dwell times with distribution functions $G_M$ and
$G_F$, with $G_F(t)\geq G_M(t)$ for all $t$ (honest fast accounts are at least as fast). For any rule that flags
accounts with dwell time below a threshold $t$, the expected number of flagged mules divided by the expected number of
flagged accounts is at most $n_M/(n_M+n_F)$.
\end{lemma}

In the simulated economy of \cref{sec:sim} with 2\% mules, arbitrageurs and honest dealers make the bound about 0.25;
the speed-only score attains a precision of 0.12. Speed is a useful input to a score, not a test.

\subsection{Exposure to flagged addresses: direction matters}\label{sec:direction}

Compliance systems commonly flag an account if the destination of its withdrawals lies within $k$ hops of a sanctioned
or frozen address. Let $G=(V,E)$ be the directed graph of transfers, with edges pointing in the direction of funds, and
let $F\subset V$ be the flagged set. Let $d_u(v,F)$ be the undirected distance from $v$ to $F$ and $d_\to(v,F)$ the length
of the shortest \emph{forward} path from $v$ into $F$. For an account $i$ with withdrawals $w$ of value $x_w$ to
destinations $v_w$, define the value-weighted exposure
\begin{equation}\label{eq:exposure}
X_i^{(k)}=\frac{\sum_w x_w\,\1\{d(v_w,F)\leq k\}}{\sum_w x_w},\qquad d\in\{d_u,d_\to\}.
\end{equation}

\begin{proposition}[Hubs and undirected distance]\label{prop:direction}
Let $h$ be an exchange hot wallet that receives at least one payment from a wallet $\ell$ that pays a flagged address
($\ell\to f$, $f\in F$), and suppose every forward path from $h$ into $F$ has length at least $m>2$, or none exists. Then
$d_u(h,F)\leq 2$ while $d_\to(h,F)\geq m$. For $2\le k<m$, every account that withdraws only to $h$ has $X^{(k)}=1$
under undirected distance and $X^{(k)}=0$ under forward distance.
\end{proposition}

Large exchanges receive deposits from criminals cashing out, so in real data almost every exchange hot wallet satisfies
the premise. Because Thai residents who bridge to offshore venues withdraw to exactly these wallets, the undirected rule
generates false positives on the population that is largest and least suspicious. Forward distance removes this
artefact but not \emph{taint contagion}: an honest remitter who pays a peer-to-peer merchant that later pays a layering
wallet is exposed under both definitions, which motivates value weighting and the combination with behavioural signals.

\subsection{Displacement across channels}\label{sec:displacement}

Why might a crackdown leave no trace in chain-wide data? Suppose a laundering network must move a fixed volume $M$ of
proceeds through channels $j\in J$---for example licensed-exchange mule accounts ($E$), peer-to-peer desks ($P$) and
over-the-counter underground bankers ($U$)---with strictly convex costs $c_j(q_j)$. Enforcement raises the marginal
cost of channel $E$ by $\tau$. The network solves
\begin{equation}\label{eq:channels}
\min_{\{q_j\geq0\}}\ \sum_{j\in J}c_j(q_j)+\tau q_E\quad\text{subject to}\quad\sum_{j\in J}q_j=M .
\end{equation}

\begin{proposition}[Displacement]\label{prop:displacement}
At an interior optimum, $\partial q_E/\partial\tau=-(1-\omega_E)/c_E''<0$ and $\partial q_j/\partial\tau=\omega_E/c_j''>0$
for $j\neq E$, where $\omega_E=(1/c_E'')/\sum_{l\in J}(1/c_l'')\in(0,1)$. Total flow is unchanged:
$\sum_j\partial q_j/\partial\tau=0$.
\end{proposition}

\begin{corollary}[Invisible crackdowns]\label{cor:invisible}
If all channels terminate in the same on-chain asset, an enforcement shock that leaves $M$ unchanged has no first-order
effect on chain-level volume. If $M$ responds through deterrence, the effect on chain-level volume is bounded by the
network's share of that volume.
\end{corollary}

\Cref{cor:invisible} is the theoretical counterpart of the event-study result in \cref{sec:eventstudy}. It also says
where to look instead: the composition of flows---the share of fast, structured deposits on licensed platforms, the
peer-to-peer premium, the network mix---moves even when the total does not.

\subsection{Typology and control points}\label{sec:typology}

\Cref{fig:typology} summarises, at the level of generality used by the FATF and UNODC, the path of fraud proceeds and
the points at which Thai institutions can intervene. The conversion step is the boundary: before it, funds sit in the
domestic banking system where they can be frozen quickly; after it, they move on a public ledger where they can be
traced indefinitely and, for USDT, frozen by the issuer.

\begin{figure}[!htbp]
\centering
\begin{tikzpicture}[
  font=\sffamily\scriptsize, >={Stealth[length=2.2mm]}, every node/.append style={execute at begin node={\hyphenpenalty=10000\exhyphenpenalty=10000}},
  stg/.style={draw=faint, line width=0.6pt, rounded corners=2pt, text width=3.4cm, minimum height=2.0cm, align=left, inner sep=5pt, fill=white},
  key/.style={draw=thai, line width=1.1pt, rounded corners=2pt, text width=3.4cm, minimum height=2.0cm, align=left, inner sep=5pt, fill=thai!8},
  ctl/.style={rounded corners=2pt, text width=3.4cm, minimum height=1.25cm, align=left, inner sep=5pt, fill=rule!70}]
\node[stg] (s1) at (0,0) {\textbf{1\; Placement}\\[2pt]Victims transfer baht to mule bank accounts after phone, chat or fake-investment
fraud; funds hop across accounts.};
\node[key, right=0.4cm of s1] (s2) {\textbf{2\; Conversion}\\[2pt]Baht become USDT via rented exchange accounts, P2P ``triangle'' trades with
unwitting sellers, or OTC desks.};
\node[stg, right=0.4cm of s2] (s3) {\textbf{3\; Layering}\\[2pt]USDT hops across Tron wallets, guarantee marketplaces, gambling sites, swaps
and bridges.};
\node[stg, right=0.4cm of s3] (s4) {\textbf{4\; Integration}\\[2pt]Exchange for yuan or dollars through USDT dealers and casinos; property and
front businesses via nominees.};
\foreach \u/\v in {s1/s2,s2/s3,s3/s4} \draw[->, draw=muted, line width=0.7pt] (\u.east) -- (\v.west);
\node[ctl, below=0.5cm of s1] (c1) {\textbf{Banks, BOT, police}\\ suspend and freeze; fraud registry; 1441 hotline};
\node[ctl, below=0.5cm of s2] (c2) {\textbf{Exchanges and SEC}\\ KYC, travel rule, address screening, mule-account bans};
\node[ctl, below=0.5cm of s3] (c3) {\textbf{Issuer and sanctions}\\ USDT freezes; OFAC listings; analytics};
\node[ctl, below=0.5cm of s4] (c4) {\textbf{AMLO, DSI, partners}\\ seizure and forfeiture; mutual legal assistance};
\foreach \s/\c in {s1/c1,s2/c2,s3/c3,s4/c4} \draw[draw=faint, line width=0.6pt, dash pattern=on 2pt off 2pt] (\s.south) -- (\c.north);
\draw[line width=0.7pt, draw=ethblue] ([yshift=0.35cm]s1.north west) -- ([yshift=0.35cm]s1.north east)
  node[midway, above, text=ethblue, font=\sffamily\tiny\bfseries] {inside the Thai banking system};
\draw[line width=0.7pt, draw=teal] ([yshift=0.35cm]s3.north west) -- ([yshift=0.35cm]s4.north east)
  node[midway, above, text=teal, font=\sffamily\tiny\bfseries] {on a public ledger: traceable, freezable by the issuer};
\draw[line width=0.7pt, draw=thai] ([yshift=0.35cm]s2.north west) -- ([yshift=0.35cm]s2.north east)
  node[midway, above, text=thai, font=\sffamily\tiny\bfseries] {the boundary};
\end{tikzpicture}
\caption{Laundering typology for fraud proceeds and the corresponding control points, at the level of generality of
FATF and UNODC typologies. The description is deliberately non-operational. Stage 2 is where funds leave the banking
system and where data held by Thai exchanges can observe them.}
\label{fig:typology}
\end{figure}
